% CPE333 Operating Systems - Lab 12: File System and Implementation
% Standalone report source; compile with pdfLaTeX or XeLaTeX.
\documentclass[a4paper,11pt]{article}
\usepackage[T1]{fontenc}
\usepackage[utf8]{inputenc}
\usepackage{lmodern}
\usepackage[margin=2.1cm]{geometry}
\usepackage{booktabs}
\usepackage{array}
\usepackage{xcolor}
\usepackage{listings}
\usepackage{enumitem}
\usepackage{needspace}
\usepackage[colorlinks=true,linkcolor=blue]{hyperref}

\definecolor{codebg}{rgb}{0.96,0.96,0.96}
\definecolor{codeframe}{rgb}{0.78,0.78,0.78}
\lstset{
  basicstyle=\ttfamily\small,
  backgroundcolor=\color{codebg},
  frame=single,
  rulecolor=\color{codeframe},
  breaklines=true,
  columns=fullflexible,
  keepspaces=true,
  showstringspaces=false,
  aboveskip=5pt,
  belowskip=5pt
}
\setlist[itemize]{itemsep=2pt,topsep=3pt}
\setlist[enumerate]{itemsep=2pt,topsep=3pt}

\title{\textbf{Lab 12 Report}\\[3pt]
       \large File System and Implementation}
\author{
Thanaboon Tikaew \quad 67070501021\\
Chanon Lhumsa-ard \quad 67070501059\\
Siriwan Yindeepot \quad 67070501075
}
\date{23 November 2026}

\begin{document}
\maketitle
\tableofcontents
\newpage

\section{Introduction}
This lab examines three file-system topics: hard and symbolic links, named pipes
(FIFOs), and the system calls used when a process reads and writes a file. The
experiments use Linux commands in WSL and Windows commands in Command Prompt
where required.

\section{Environment}
\begin{tabular}{>{\bfseries}p{0.28\linewidth}p{0.66\linewidth}}
Linux environment & WSL 2; Ubuntu 24.04.1 LTS (noble), from \texttt{lsb\_release -a} \\
Linux working directory & \texttt{\textasciitilde/CPE333/PS12} \\
Windows environment & Windows 11 Pro, version 10.0.26300 (build 26300) \\
Windows shell & Command Prompt (\texttt{cmd.exe}) for \texttt{mklink} \\
\end{tabular}

\noindent\textbf{Note:} The experiments were performed by different group
members on their own systems. Usernames and available drive letters in the
captured output therefore vary. The cross-drive run was performed on a Windows
system with a separate D: volume labeled \texttt{DATA}; some members' systems
did not have a D: drive.

\section{Hard and Symbolic Links}

Together these links turn the directory hierarchy from a strict tree into a
general graph directory structure: a hard link gives one file more than one
parent (a DAG), while a symbolic link stores a pathname and can also form a
cycle when it points back to an upper-level directory.

\subsection{Create and inspect links}
In Linux, a hard link is another directory entry for the same file data, while
a symbolic link stores a path to its target. The inode numbers shown by
\texttt{ls -li} help distinguish the two.

\begin{lstlisting}
# Linux (WSL)
mkdir -p ~/CPE333/PS12/link-lab
cd ~/CPE333/PS12/link-lab
printf 'hello\n' > source.txt
ln source.txt hard.txt
ln -s source.txt soft.txt
ls -li source.txt hard.txt soft.txt
\end{lstlisting}

On Windows, run the following in Command Prompt in a test directory:
\begin{lstlisting}
cd /d "%USERPROFILE%"
if not exist PS12-link-lab mkdir PS12-link-lab
cd PS12-link-lab
echo hello>source.txt
mklink /H hard.txt source.txt
mklink soft.txt source.txt
dir
\end{lstlisting}

\noindent Creating Windows symbolic links may require Developer Mode or an
elevated Command Prompt, depending on system policy.

\subsection{Delete the source}
After recording the directory listing, remove only the test source file and
read both links:
\begin{lstlisting}
# Linux (WSL)
rm source.txt
cat hard.txt
cat soft.txt
\end{lstlisting}
\begin{lstlisting}
:: Windows Command Prompt
del source.txt
type hard.txt
type soft.txt
\end{lstlisting}

\begin{itemize}
  \item Expected behavior: the hard link still refers to the file data because
  it is another directory entry for that file. The symbolic link refers to the
  old pathname, so it becomes dangling after the target is removed.
\end{itemize}

\noindent\textbf{Observed output from WSL:}
\begin{lstlisting}
1645 -rw-r--r-- 2 thnabun thnabun  6 Nov 21 18:09 hard.txt
12454 lrwxrwxrwx 1 thnabun thnabun 10 Nov 21 18:09 soft.txt -> source.txt
1645 -rw-r--r-- 2 thnabun thnabun  6 Nov 21 18:09 source.txt

$ rm source.txt
$ cat hard.txt
hello
$ cat soft.txt
cat: soft.txt: No such file or directory
\end{lstlisting}

The matching inode number \texttt{1645} and link count of 2 show that
\texttt{source.txt} and \texttt{hard.txt} were two names for the same file
data. Removing \texttt{source.txt} left \texttt{hard.txt} readable. The
symbolic link stored the relative pathname \texttt{source.txt}; after that
pathname was removed, \texttt{soft.txt} could no longer resolve its target.

\noindent\textbf{Observed Windows result:}
\begin{lstlisting}
mklink /H hard.txt source.txt
Hardlink created for hard.txt <<===>> source.txt
mklink soft.txt source.txt
You do not have sufficient privilege to perform this operation.

:: Repeated in an elevated Command Prompt
mklink soft.txt source.txt
symbolic link created for soft.txt <<===>> source.txt
dir
<SYMLINK> soft.txt [source.txt]
del source.txt
type hard.txt
hello
type soft.txt
The system cannot find the file specified.
\end{lstlisting}

The Windows hard link was created successfully and remained readable after
\texttt{source.txt} was deleted. The first symbolic-link attempt was denied
for insufficient privilege. Repeating it in an elevated Command Prompt
created \texttt{soft.txt}, which \texttt{dir} identified as a symbolic link to
\texttt{source.txt}. After deleting the target, \texttt{type soft.txt} could
not open it, confirming the dangling-link behavior.

\subsection{A two-level symbolic-link chain}
Recreate the source in the test directory, then create a link to the source
and a second link to the first link:
\begin{lstlisting}
# Linux (WSL)
printf 'chain content\n' > source.txt
ln -sfn source.txt level1
ln -sfn level1 level2
cat level2
rm source.txt
cat level2
\end{lstlisting}
\begin{lstlisting}
:: Windows Command Prompt
echo chain content>source.txt
mklink level1 source.txt
mklink level2 level1
type level2
del source.txt
type level2
\end{lstlisting}

Then recreate the source and links, remove only \texttt{level1}, and read
\texttt{level2}. The second link depends on resolving the first link, so
removing the first link can make the second link unusable even while the
original source file exists.

\begin{lstlisting}
# Linux (WSL), reset the chain for the second deletion test
printf 'chain content\n' > source.txt
ln -sfn source.txt level1
ln -sfn level1 level2
rm level1
cat level2
\end{lstlisting}
\begin{lstlisting}
:: Windows Command Prompt, reset the chain for the second deletion test
del level2
del level1
echo chain content>source.txt
mklink level1 source.txt
mklink level2 level1
del level1
type level2
\end{lstlisting}

\noindent\textbf{Observed output after deleting the source:}
\begin{lstlisting}
$ cat level2
chain content
$ rm source.txt
$ cat level2
cat: level2: No such file or directory
\end{lstlisting}

Before deletion, resolving \texttt{level2} followed the chain
\texttt{level2} $\to$ \texttt{level1} $\to$ \texttt{source.txt}. Removing
\texttt{source.txt} broke the final link in the chain, so \texttt{cat level2}
could no longer open the target.

\noindent\textbf{Observed Windows result after deleting the source:}
\begin{lstlisting}
> type level2
chain content

> del source.txt

> type level2
The system cannot find the file specified.
\end{lstlisting}

Windows behaved in the same way. \texttt{level2} first resolved through
\texttt{level1} to \texttt{source.txt}; after the final target was deleted,
both symbolic links still existed but the chain could no longer be resolved.

\noindent\textbf{Observed result after deleting \texttt{level1}:}
\begin{lstlisting}
$ cat level2
chain content
$ rm level1
$ cat level2
cat: level2: No such file or directory
\end{lstlisting}

At this point \texttt{source.txt} still existed, but \texttt{level2} could not
resolve its immediate target \texttt{level1}. Removing an intermediate
symbolic link therefore broke the chain.

\Needspace{8\baselineskip}
\noindent\textbf{Observed Windows result after deleting \texttt{level1}:}
\begin{lstlisting}
> type level2
chain content

> del level1

> type level2
The system cannot find the file specified.
\end{lstlisting}

The Windows result matched Linux: \texttt{source.txt} was still present, but
\texttt{level2} could not resolve its immediate target \texttt{level1}.

\subsection{Directory links and recursive listing}
Create a directory link that points from a child directory back to its parent:
\begin{lstlisting}
# Linux (WSL)
mkdir -p tree/child
ln -s tree treeLink
ln -s .. tree/child/parent
ls -R tree
\end{lstlisting}
\begin{lstlisting}
:: Windows Command Prompt
mkdir tree\child
mklink /D treeLink tree
mklink /D tree\child\parent ..
dir /s tree
\end{lstlisting}

Record whether each recursive listing follows the link, skips it, or reports
an error. Stop the Windows listing with \texttt{Ctrl+C} if it continues
repeating directories. On Linux, the default \texttt{ls -R} does not follow
symbolic links to directories.

\noindent\textbf{Observed WSL output:}
\begin{lstlisting}
$ ls -ld treeLink tree/child/parent
lrwxrwxrwx 1 thnabun thnabun 2 Nov 21 18:22 tree/child/parent -> ..
lrwxrwxrwx 1 thnabun thnabun 4 Nov 21 18:22 treeLink -> tree

$ ls -R tree
tree:
child

tree/child:
parent
\end{lstlisting}

Both directory symbolic links were created. The \texttt{parent} link points
back to \texttt{tree}, but the default \texttt{ls -R} listed the link without
following it, so the recursive listing did not enter a loop.

\noindent\textbf{Observed Windows output (excerpt):}
\begin{lstlisting}
mklink /D treeLink tree
symbolic link created for treeLink <<===>> tree
mklink /D tree\child\parent ..
symbolic link created for tree\child\parent <<===>> ..

Directory of ...\tree\child
<SYMLINKD> parent [..]
Directory of ...\tree\child\parent\child
<SYMLINKD> parent [..]
...
Total Files Listed: 0
384 Dir(s)
\end{lstlisting}

Windows \texttt{dir /s} followed the directory link into a parent-child cycle
and listed 384 directories before stopping; it did not report why.

\subsection{Links across drives}
If the computer has a D: drive, test a symbolic link on D: pointing to a
source file on C:. Use an existing target file and a new link path:
\begin{lstlisting}
:: Windows Command Prompt; replace the target path if needed
if not exist D:\PS12-link-lab mkdir D:\PS12-link-lab
echo cross-drive test>"%USERPROFILE%\PS12-link-lab\source.txt"
mklink D:\PS12-link-lab\cross-drive.txt "%USERPROFILE%\PS12-link-lab\source.txt"
mklink /H D:\PS12-link-lab\hard-cross-drive.txt "%USERPROFILE%\PS12-link-lab\source.txt"
\end{lstlisting}

The symbolic-link test may be affected by Windows permissions and policy, so
the commands were run in an elevated Command Prompt. Hard links cannot cross
volumes, so the \texttt{/H} command is expected to fail when C: and D: are
different volumes.

\noindent\textbf{Observed result (Windows, elevated Command Prompt):}
\begin{lstlisting}
mklink D:\PS12-link-lab\cross-drive.txt "%USERPROFILE%\PS12-link-lab\source.txt"
symbolic link created for D:\PS12-link-lab\cross-drive.txt <<===>> C:\Users\win25\PS12-link-lab\source.txt

mklink /H D:\PS12-link-lab\hard-cross-drive.txt "%USERPROFILE%\PS12-link-lab\source.txt"
The system cannot move the file to a different disk drive.

dir D:\PS12-link-lab
 Volume in drive D is DATA
 ...
11/23/2026  11:54 AM    <SYMLINK>      cross-drive.txt [C:\Users\win25\PS12-link-lab\source.txt]

type D:\PS12-link-lab\cross-drive.txt
cross-drive test
\end{lstlisting}

On the member's Windows PC, D: was labeled \texttt{DATA}. The D: symbolic link
read the C: source, while the hard-link attempt returned the different-drive
error because hard links must share a volume.

\section{Linux Named Pipe}

\subsection{Create and use a FIFO}
Run this in WSL: a FIFO has a directory entry but passes data between processes
instead of storing it in a regular file.

\begin{lstlisting}
# Prepare the FIFO (run once)
mkdir -p ~/CPE333/PS12/pipe-lab
cd ~/CPE333/PS12/pipe-lab
mkfifo pipeFile
ls -la pipeFile
\end{lstlisting}

Open two WSL terminals in this directory. Run the matching command in each,
then type lines in Terminal 1 and press \texttt{Ctrl+D} to close the writer.
\begin{lstlisting}
# Terminal 1 (writer)
cat > pipeFile
# Terminal 2 (reader)
cat pipeFile
\end{lstlisting}

\Needspace{14\baselineskip}
\noindent\textbf{Observed output:}
\begin{lstlisting}
$ ls -la pipeFile
prw-r--r-- 1 thnabun thnabun 0 Nov 23 11:54 pipeFile

Terminal 1: cat > pipeFile
Typed: hello linux
Typed: second line

Terminal 2: cat pipeFile
Displayed: hello linux
Displayed: second line
\end{lstlisting}

The leading \texttt{p} identifies \texttt{pipeFile} as a FIFO. Its size is zero
because it does not store data as a regular file. Opening one end waits for the
other; once both \texttt{cat} processes are connected, lines typed in Terminal 1
are sent to Terminal 2 in order. The FIFO is a unidirectional byte stream.
\texttt{Ctrl+D} closes the writer, and the reader exits after receiving
end-of-file.

\section{Reading and Writing Files}

\subsection{Trace the system calls}
The following command sequence was run in WSL from
\texttt{\textasciitilde/CPE333/PS12}:
\begin{lstlisting}
echo hello > foo
strace cat foo
\end{lstlisting}
The trace is written to standard error while \texttt{cat} writes the file
contents to standard output, so the two outputs may appear interleaved in a
terminal.

The command created a six-byte file containing \texttt{hello} followed by a
newline. The trace contained many setup calls; the four calls relevant to the
question were:
\begin{lstlisting}
read(3, "hello\n", 131072) = 6
write(1, "hello\n", 6) = 6
read(3, "", 131072) = 0
close(3) = 0
\end{lstlisting}

\subsection{Explanation}
\begin{itemize}
  \item \texttt{3} is the file descriptor assigned to \texttt{foo}. The trace
  shows the earlier \texttt{openat(..., "foo", ...) = 3}.
  \item \texttt{1} is the standard-output file descriptor. \texttt{cat} writes
  the file contents to standard output.
  \item \texttt{6} is the number of bytes actually read and written:
  \texttt{hello} has five characters and the newline adds one byte.
  \item The third argument to \texttt{read} is its requested byte count:
  \texttt{32768} in the assignment example and \texttt{131072} in this run.
  It is a maximum, not the amount actually read. Here \texttt{read} returned
  \texttt{6} bytes, then \texttt{0} at end-of-file.
\end{itemize}

\end{document}
